\section{Einleitung}

Künstliche Intelligenz (KI) hat sich zu einem der relevantesten Bereiche der modernen Informatik entwickelt.
In den letzten Jahren haben der Fortschritt der Rechenleistung, die Verfügbarkeit von Daten und die Modelle
des maschinellen Lernens die Schaffung von Systemen ermöglicht, die Aufgaben ausführen können, die
zuvor ausschließlich von menschlicher Intervention abhingen.

Derzeit sind KI-Anwendungen in verschiedenen Bereichen der Gesellschaft präsent, darunter Gesundheit,
Industrie, Transport, Bildung, Sicherheit und Unterhaltung. Empfehlungssysteme,
virtuelle Assistenten, automatische Übersetzung und autonome Fahrzeuge sind nur einige Beispiele
weit verbreiteter Technologien.

Neben dem technologischen Einfluss beeinflusst KI auch wirtschaftliche und soziale Aspekte. Unternehmen
investieren kontinuierlich in Automatisierung und Datenanalyse, um die Produktivität zu verbessern,
Kosten zu senken und die Wettbewerbsfähigkeit auf dem globalen Markt zu steigern. 